\section{Background}\label{sec:background}

Globally, e-commerce (or online shopping) has been a significant driver of revenue for business as well as a vital utility for consumers to access goods and services that may have otherwise been difficult to access.
Total e-commerce sales across the globe was projected to surpass 6.86 trillion USD by the end of 2025\cite{sellerscommerce2026ecommerce}.
While western markets such as North America and Europe have traditionally dominated this sector, the focus has been rapidly shifting towards Asia.
South Asia in particular has been a region of strong growth with a Compound Annual Growth Rate (CAGR) exceeding 20\%\cite{koutsou-wehling2025southeast}.

Sri Lanka is a notable player in this rapidly growing subspace, with its digital economy being valued at 1,342 billion LKR in 2025, contributing 4.5\% to the national GDP\cite{dailyft2025lankan}.
This growth was accelerated during the COVID-19 pandemic, which forced a migration of both consumers and SMEs to online shopping.
Many of these e-commerce platforms offer consumers access to products that are otherwise unavailable in their localities, while also expanding the reach for domestic businesses\cite{slasscom2023daraz}.
However, a barrier for growth, is the lack of consumer trust in online shopping which is currently plagued by a “fake discount epidemic”, with anti-consumer practices such as false reference pricing, perpetual sales, and bait and switch techniques \cite{devarmani2025fake}.

Data from neighbouring countries such as India have shown that over 50\% of consumers who engage in online shopping have fallen victim to such scams\cite{economictimes2019over}.
This pattern seems to be mirrored in Sri Lanka as well, where profit maximization is often prioritized over consumer satisfaction.
Moreover, the lack of a centralized data source to compare product features and prices also prevents finding the best deal, often a very labour intensive and time-consuming process for the average user.
Modern consumers are also aware of “dynamic pricing”, which is the use of algorithms to dynamically adjust prices of products based on internal factors such as customer profiling as well as external trends such as demand and availability.
Online shoppers are driven to monitor product prices and wait for markdowns before making a purchase\cite{victor2018price}.
From a data perspective, Sri Lanka’s e-commerce landscape is highly fragmented.

This project aims to bridge this information gap by performing a comprehensive analysis of Sri Lanka’s local e-commerce ecosystem, centered around consumer goods (which is the largest single category in the nation, with over 76\% of e-commerce users have purchased consumer goods online at least once\cite{slasscom2023daraz})),  By leveraging techniques of big data analytics, we will also aim to identify pricing correlations with external events, understand the prevalence of the “fake discount epidemic”, and develop a generalized model which is capable of automated data extractions.
Ultimately, this work will facilitate the creation of tools for price tracking and cross-site product comparison, empowering consumers and providing the data-driven insights necessary to foster innovation and transparency in Sri Lanka’s digital economy.
